\documentclass[11pt]{article}
% ---------------------------------------------------------------
% R3.07 SQL dans un langage de programmation -- preambule commun des sujets de TP
% Article clair, fond blanc, accents de couleur coherents avec les slides
% Charte derivee de R5.A.08 Qualite de developpement (meme palette IUT)
% Utilisation :
%   \documentclass[11pt]{article}
%   % ---------------------------------------------------------------
% R3.07 SQL dans un langage de programmation -- preambule commun des sujets de TP
% Article clair, fond blanc, accents de couleur coherents avec les slides
% Charte derivee de R5.A.08 Qualite de developpement (meme palette IUT)
% Utilisation :
%   \documentclass[11pt]{article}
%   % ---------------------------------------------------------------
% R3.07 SQL dans un langage de programmation -- preambule commun des sujets de TP
% Article clair, fond blanc, accents de couleur coherents avec les slides
% Charte derivee de R5.A.08 Qualite de developpement (meme palette IUT)
% Utilisation :
%   \documentclass[11pt]{article}
%   \input{../../latex/common/r307-tp-preamble.tex}
%   \begin{document}
%   \rtpheader{1}{Titre du TP}
% ---------------------------------------------------------------
\usepackage[T1]{fontenc}
\usepackage[utf8]{inputenc}
\usepackage[french]{babel}
\usepackage[scale=0.92]{tgheros}
\usepackage{tgpagella}
\usepackage[a4paper, margin=2.2cm, headheight=14pt]{geometry}
\usepackage{graphicx}
\usepackage{xcolor}
\usepackage{listings}
\usepackage[skins,listings,breakable]{tcolorbox}
\usepackage{enumitem}
\usepackage{booktabs}
\usepackage{array}
\usepackage{multirow}
\usepackage{tikz}
\usetikzlibrary{shapes, shapes.geometric, arrows, arrows.meta, positioning, calc}
\usepackage[colorlinks=true, linkcolor=qaaccentdark, urlcolor=qaaccentdark]{hyperref}
\usepackage{fancyhdr}
\usepackage{lastpage}
\usepackage{titlesec}
\usepackage{etoolbox}

% ---------------------------------------------------------------
% PALETTE -- version imprimable des couleurs des diapositives
% ---------------------------------------------------------------
\definecolor{qaaccentdark}{HTML}{7A1680}   % violet du logo IUT, fonce pour rester lisible sur blanc
\definecolor{qaaccent2dark}{HTML}{9B4B94}  % violet clair du logo IUT (accent secondaire)
\definecolor{qapanellight}{HTML}{F4F6F8}   % fond de bloc tres clair
\definecolor{qalinelight}{HTML}{D8DEE4}    % filets
\definecolor{qatextdark}{HTML}{1A1D23}     % texte principal
\definecolor{qamuteddark}{HTML}{5B6472}    % texte secondaire
\definecolor{qacodekwl}{HTML}{7A1680}      % code : mots-cles
\definecolor{qacodestrl}{HTML}{2E7D32}     % code : chaines
\definecolor{qacodecoml}{HTML}{6B7280}     % code : commentaires
\definecolor{qacodeannl}{HTML}{B26A00}     % code : annotations / javadoc

\color{qatextdark}

\setlength{\parskip}{4pt}
\setlength{\parindent}{0pt}

% ---------------------------------------------------------------
% TITRES DE SECTION -- texte blanc sur bande de couleur pleine largeur
% (gere aussi les variantes etoilees \section*)
% ---------------------------------------------------------------
\newcommand{\qasectionband}[2]{%
  \noindent\colorbox{#1}{%
    \parbox[c]{\dimexpr\linewidth-2\fboxsep\relax}{%
      \color{white}#2%
    }%
  }%
}
\newcommand{\secformat}[1]{\qasectionband{qaaccentdark}{\bfseries\Large\thesection\quad #1}}
\newcommand{\secformatnn}[1]{\qasectionband{qaaccentdark}{\bfseries\Large #1}}
\newcommand{\subsecformat}[1]{\qasectionband{qaaccent2dark}{\bfseries\large\thesubsection\quad #1}}
\newcommand{\subsecformatnn}[1]{\qasectionband{qaaccent2dark}{\bfseries\large #1}}

\titleformat{\section}{\normalfont}{}{0pt}{\secformat}
\titleformat{name=\section,numberless}{\normalfont}{}{0pt}{\secformatnn}
\titleformat{\subsection}{\normalfont}{}{0pt}{\subsecformat}
\titleformat{name=\subsection,numberless}{\normalfont}{}{0pt}{\subsecformatnn}
\titleformat{\subsubsection}{\normalfont\bfseries\color{qaaccentdark}}{\thesubsubsection}{0.6em}{}
\titlespacing*{\section}{0pt}{14pt}{8pt}
\titlespacing*{\subsection}{0pt}{10pt}{6pt}
\titlespacing*{\subsubsection}{0pt}{8pt}{4pt}

% ---------------------------------------------------------------
% LISTINGS -- theme clair (lstlisting reste utilisable tel quel)
% \lstcode : police + couleur de base, a reutiliser dans les surcharges
% ---------------------------------------------------------------
\newcommand{\lstcode}{\ttfamily\color{qatextdark}}

\lstdefinelanguage{TypeScript}{
  keywords={typeof, new, true, false, catch, function, return, null, switch, var,
    if, in, while, do, else, case, break, const, let, interface, class, extends,
    implements, enum, export, import, this, super, constructor, public, private,
    protected, static, async, await, throw, throws, try, finally, void},
  ndkeywords={boolean, any, number, string, int, long, double},
  sensitive=true,
  comment=[l]{//},
  morecomment=[s]{/*}{*/},
  morestring=[b]',
  morestring=[b]"
}

\lstdefinestyle{r307code}{
  basicstyle=\lstcode\footnotesize,
  keywordstyle=\color{qacodekwl}\bfseries,
  ndkeywordstyle=\color{qaaccent2dark}\bfseries,
  stringstyle=\color{qacodestrl},
  commentstyle=\color{qacodecoml}\itshape,
  morecomment=[s][\color{qacodeannl}]{/**}{*/},
  backgroundcolor=\color{qapanellight},
  rulecolor=\color{qalinelight},
  frame=single,
  framesep=4pt,
  xleftmargin=6pt,
  xrightmargin=6pt,
  numbers=left,
  numberstyle=\tiny\color{qamuteddark},
  numbersep=8pt,
  keepspaces=true,
  columns=fullflexible,
  showstringspaces=false,
  showspaces=false,
  showtabs=false,
  breaklines=true,
  breakatwhitespace=false,
  captionpos=b,
  tabsize=2,
  escapeinside={(*@}{@*)},
}
\lstset{style=r307code}

% ---------------------------------------------------------------
% BOITES : exercice, code, information
% ---------------------------------------------------------------
\tcbset{
  exobox/.style={colback=qapanellight, colframe=qaaccentdark, boxrule=0.6pt,
    arc=2pt, left=8pt, right=8pt, top=6pt, bottom=6pt,
    coltitle=white, colbacktitle=qaaccentdark, fonttitle=\bfseries,
    boxsep=2pt, breakable},
  infobox/.style={colback=white, colframe=qaaccent2dark, boxrule=0.6pt,
    leftrule=3pt, arc=0pt, left=8pt, right=8pt, top=5pt, bottom=5pt,
    coltitle=qaaccent2dark, colbacktitle=white, fonttitle=\bfseries,
    breakable},
  codebox/.style={colback=qatextdark, colframe=qalinelight, boxrule=0.4pt,
    arc=2pt, left=8pt, right=8pt, top=5pt, bottom=5pt, listing only,
    listing options={basicstyle=\ttfamily\small\color{white},
      breaklines=true, showstringspaces=false, columns=fullflexible,
      keepspaces=true, frame=none, numbers=none, backgroundcolor={}}, breakable}
}
\newtcolorbox[auto counter]{exobox}[2][]{exobox,
  title={Exercice~\thetcbcounter\ifstrempty{#2}{}{ --- #2}}, #1}
\newtcolorbox{infobox}[1][\`A retenir]{infobox, title={#1}}
\newtcblisting{codebox}{codebox}

% ---------------------------------------------------------------
% EN-TETE DE PAGE (fil d'Ariane + pagination)
% ---------------------------------------------------------------
\newcommand{\qaheadtitlevalue}{}
\newcommand{\qaheadtitle}[1]{\renewcommand{\qaheadtitlevalue}{#1}}

\pagestyle{fancy}
\fancyhf{}
\renewcommand{\headrulewidth}{0.4pt}
\fancyhead[L]{\small\color{qamuteddark}R3.07 --- \qaheadtitlevalue}
\fancyhead[R]{\small\color{qamuteddark}\thepage\,/\,\pageref{LastPage}}
\fancyfoot[C]{\footnotesize\color{qamuteddark}IUT de Vannes --- Universit\'e Bretagne Sud}
\fancypagestyle{plain}{\fancyhf{}\renewcommand{\headrulewidth}{0pt}%
  \fancyfoot[C]{\footnotesize\color{qamuteddark}IUT de Vannes --- Universit\'e Bretagne Sud}}

% ---------------------------------------------------------------
% BANDEAU DE TITRE (pas une page dediee, pour rester compact)
% \rtpheader{numero ou intitule court}{titre du TP}
% ---------------------------------------------------------------
\newcommand{\ryear}{2026/2027}
\newcommand{\rtpheader}[2]{%
  \thispagestyle{plain}%
  \begin{minipage}[t]{0.68\textwidth}
    {\small\color{qamuteddark}
      Universit\'e Bretagne Sud\\
      IUT de Vannes\\
      D\'epartement Informatique --- \ryear}
  \end{minipage}%
  \hfill
  \begin{minipage}[t]{0.28\textwidth}
    \raggedleft
    \includegraphics[width=2.6cm]{iut_logo.jpg}
  \end{minipage}

  \vspace{6mm}
  \begin{center}
    {\large\bfseries\color{qatextdark} BUT 2 --- R3.07. SQL dans un langage de programmation}\\[2mm]
    {\normalsize\color{qamuteddark} TD/TP #1}\\[1mm]
    {\Large\bfseries\color{qaaccentdark} #2}
  \end{center}
  \vspace{4mm}
  \qaheadtitle{#2}
}

%   \begin{document}
%   \rtpheader{1}{Titre du TP}
% ---------------------------------------------------------------
\usepackage[T1]{fontenc}
\usepackage[utf8]{inputenc}
\usepackage[french]{babel}
\usepackage[scale=0.92]{tgheros}
\usepackage{tgpagella}
\usepackage[a4paper, margin=2.2cm, headheight=14pt]{geometry}
\usepackage{graphicx}
\usepackage{xcolor}
\usepackage{listings}
\usepackage[skins,listings,breakable]{tcolorbox}
\usepackage{enumitem}
\usepackage{booktabs}
\usepackage{array}
\usepackage{multirow}
\usepackage{tikz}
\usetikzlibrary{shapes, shapes.geometric, arrows, arrows.meta, positioning, calc}
\usepackage[colorlinks=true, linkcolor=qaaccentdark, urlcolor=qaaccentdark]{hyperref}
\usepackage{fancyhdr}
\usepackage{lastpage}
\usepackage{titlesec}
\usepackage{etoolbox}

% ---------------------------------------------------------------
% PALETTE -- version imprimable des couleurs des diapositives
% ---------------------------------------------------------------
\definecolor{qaaccentdark}{HTML}{7A1680}   % violet du logo IUT, fonce pour rester lisible sur blanc
\definecolor{qaaccent2dark}{HTML}{9B4B94}  % violet clair du logo IUT (accent secondaire)
\definecolor{qapanellight}{HTML}{F4F6F8}   % fond de bloc tres clair
\definecolor{qalinelight}{HTML}{D8DEE4}    % filets
\definecolor{qatextdark}{HTML}{1A1D23}     % texte principal
\definecolor{qamuteddark}{HTML}{5B6472}    % texte secondaire
\definecolor{qacodekwl}{HTML}{7A1680}      % code : mots-cles
\definecolor{qacodestrl}{HTML}{2E7D32}     % code : chaines
\definecolor{qacodecoml}{HTML}{6B7280}     % code : commentaires
\definecolor{qacodeannl}{HTML}{B26A00}     % code : annotations / javadoc

\color{qatextdark}

\setlength{\parskip}{4pt}
\setlength{\parindent}{0pt}

% ---------------------------------------------------------------
% TITRES DE SECTION -- texte blanc sur bande de couleur pleine largeur
% (gere aussi les variantes etoilees \section*)
% ---------------------------------------------------------------
\newcommand{\qasectionband}[2]{%
  \noindent\colorbox{#1}{%
    \parbox[c]{\dimexpr\linewidth-2\fboxsep\relax}{%
      \color{white}#2%
    }%
  }%
}
\newcommand{\secformat}[1]{\qasectionband{qaaccentdark}{\bfseries\Large\thesection\quad #1}}
\newcommand{\secformatnn}[1]{\qasectionband{qaaccentdark}{\bfseries\Large #1}}
\newcommand{\subsecformat}[1]{\qasectionband{qaaccent2dark}{\bfseries\large\thesubsection\quad #1}}
\newcommand{\subsecformatnn}[1]{\qasectionband{qaaccent2dark}{\bfseries\large #1}}

\titleformat{\section}{\normalfont}{}{0pt}{\secformat}
\titleformat{name=\section,numberless}{\normalfont}{}{0pt}{\secformatnn}
\titleformat{\subsection}{\normalfont}{}{0pt}{\subsecformat}
\titleformat{name=\subsection,numberless}{\normalfont}{}{0pt}{\subsecformatnn}
\titleformat{\subsubsection}{\normalfont\bfseries\color{qaaccentdark}}{\thesubsubsection}{0.6em}{}
\titlespacing*{\section}{0pt}{14pt}{8pt}
\titlespacing*{\subsection}{0pt}{10pt}{6pt}
\titlespacing*{\subsubsection}{0pt}{8pt}{4pt}

% ---------------------------------------------------------------
% LISTINGS -- theme clair (lstlisting reste utilisable tel quel)
% \lstcode : police + couleur de base, a reutiliser dans les surcharges
% ---------------------------------------------------------------
\newcommand{\lstcode}{\ttfamily\color{qatextdark}}

\lstdefinelanguage{TypeScript}{
  keywords={typeof, new, true, false, catch, function, return, null, switch, var,
    if, in, while, do, else, case, break, const, let, interface, class, extends,
    implements, enum, export, import, this, super, constructor, public, private,
    protected, static, async, await, throw, throws, try, finally, void},
  ndkeywords={boolean, any, number, string, int, long, double},
  sensitive=true,
  comment=[l]{//},
  morecomment=[s]{/*}{*/},
  morestring=[b]',
  morestring=[b]"
}

\lstdefinestyle{r307code}{
  basicstyle=\lstcode\footnotesize,
  keywordstyle=\color{qacodekwl}\bfseries,
  ndkeywordstyle=\color{qaaccent2dark}\bfseries,
  stringstyle=\color{qacodestrl},
  commentstyle=\color{qacodecoml}\itshape,
  morecomment=[s][\color{qacodeannl}]{/**}{*/},
  backgroundcolor=\color{qapanellight},
  rulecolor=\color{qalinelight},
  frame=single,
  framesep=4pt,
  xleftmargin=6pt,
  xrightmargin=6pt,
  numbers=left,
  numberstyle=\tiny\color{qamuteddark},
  numbersep=8pt,
  keepspaces=true,
  columns=fullflexible,
  showstringspaces=false,
  showspaces=false,
  showtabs=false,
  breaklines=true,
  breakatwhitespace=false,
  captionpos=b,
  tabsize=2,
  escapeinside={(*@}{@*)},
}
\lstset{style=r307code}

% ---------------------------------------------------------------
% BOITES : exercice, code, information
% ---------------------------------------------------------------
\tcbset{
  exobox/.style={colback=qapanellight, colframe=qaaccentdark, boxrule=0.6pt,
    arc=2pt, left=8pt, right=8pt, top=6pt, bottom=6pt,
    coltitle=white, colbacktitle=qaaccentdark, fonttitle=\bfseries,
    boxsep=2pt, breakable},
  infobox/.style={colback=white, colframe=qaaccent2dark, boxrule=0.6pt,
    leftrule=3pt, arc=0pt, left=8pt, right=8pt, top=5pt, bottom=5pt,
    coltitle=qaaccent2dark, colbacktitle=white, fonttitle=\bfseries,
    breakable},
  codebox/.style={colback=qatextdark, colframe=qalinelight, boxrule=0.4pt,
    arc=2pt, left=8pt, right=8pt, top=5pt, bottom=5pt, listing only,
    listing options={basicstyle=\ttfamily\small\color{white},
      breaklines=true, showstringspaces=false, columns=fullflexible,
      keepspaces=true, frame=none, numbers=none, backgroundcolor={}}, breakable}
}
\newtcolorbox[auto counter]{exobox}[2][]{exobox,
  title={Exercice~\thetcbcounter\ifstrempty{#2}{}{ --- #2}}, #1}
\newtcolorbox{infobox}[1][\`A retenir]{infobox, title={#1}}
\newtcblisting{codebox}{codebox}

% ---------------------------------------------------------------
% EN-TETE DE PAGE (fil d'Ariane + pagination)
% ---------------------------------------------------------------
\newcommand{\qaheadtitlevalue}{}
\newcommand{\qaheadtitle}[1]{\renewcommand{\qaheadtitlevalue}{#1}}

\pagestyle{fancy}
\fancyhf{}
\renewcommand{\headrulewidth}{0.4pt}
\fancyhead[L]{\small\color{qamuteddark}R3.07 --- \qaheadtitlevalue}
\fancyhead[R]{\small\color{qamuteddark}\thepage\,/\,\pageref{LastPage}}
\fancyfoot[C]{\footnotesize\color{qamuteddark}IUT de Vannes --- Universit\'e Bretagne Sud}
\fancypagestyle{plain}{\fancyhf{}\renewcommand{\headrulewidth}{0pt}%
  \fancyfoot[C]{\footnotesize\color{qamuteddark}IUT de Vannes --- Universit\'e Bretagne Sud}}

% ---------------------------------------------------------------
% BANDEAU DE TITRE (pas une page dediee, pour rester compact)
% \rtpheader{numero ou intitule court}{titre du TP}
% ---------------------------------------------------------------
\newcommand{\ryear}{2026/2027}
\newcommand{\rtpheader}[2]{%
  \thispagestyle{plain}%
  \begin{minipage}[t]{0.68\textwidth}
    {\small\color{qamuteddark}
      Universit\'e Bretagne Sud\\
      IUT de Vannes\\
      D\'epartement Informatique --- \ryear}
  \end{minipage}%
  \hfill
  \begin{minipage}[t]{0.28\textwidth}
    \raggedleft
    \includegraphics[width=2.6cm]{iut_logo.jpg}
  \end{minipage}

  \vspace{6mm}
  \begin{center}
    {\large\bfseries\color{qatextdark} BUT 2 --- R3.07. SQL dans un langage de programmation}\\[2mm]
    {\normalsize\color{qamuteddark} TD/TP #1}\\[1mm]
    {\Large\bfseries\color{qaaccentdark} #2}
  \end{center}
  \vspace{4mm}
  \qaheadtitle{#2}
}

%   \begin{document}
%   \rtpheader{1}{Titre du TP}
% ---------------------------------------------------------------
\usepackage[T1]{fontenc}
\usepackage[utf8]{inputenc}
\usepackage[french]{babel}
\usepackage[scale=0.92]{tgheros}
\usepackage{tgpagella}
\usepackage[a4paper, margin=2.2cm, headheight=14pt]{geometry}
\usepackage{graphicx}
\usepackage{xcolor}
\usepackage{listings}
\usepackage[skins,listings,breakable]{tcolorbox}
\usepackage{enumitem}
\usepackage{booktabs}
\usepackage{array}
\usepackage{multirow}
\usepackage{tikz}
\usetikzlibrary{shapes, shapes.geometric, arrows, arrows.meta, positioning, calc}
\usepackage[colorlinks=true, linkcolor=qaaccentdark, urlcolor=qaaccentdark]{hyperref}
\usepackage{fancyhdr}
\usepackage{lastpage}
\usepackage{titlesec}
\usepackage{etoolbox}

% ---------------------------------------------------------------
% PALETTE -- version imprimable des couleurs des diapositives
% ---------------------------------------------------------------
\definecolor{qaaccentdark}{HTML}{7A1680}   % violet du logo IUT, fonce pour rester lisible sur blanc
\definecolor{qaaccent2dark}{HTML}{9B4B94}  % violet clair du logo IUT (accent secondaire)
\definecolor{qapanellight}{HTML}{F4F6F8}   % fond de bloc tres clair
\definecolor{qalinelight}{HTML}{D8DEE4}    % filets
\definecolor{qatextdark}{HTML}{1A1D23}     % texte principal
\definecolor{qamuteddark}{HTML}{5B6472}    % texte secondaire
\definecolor{qacodekwl}{HTML}{7A1680}      % code : mots-cles
\definecolor{qacodestrl}{HTML}{2E7D32}     % code : chaines
\definecolor{qacodecoml}{HTML}{6B7280}     % code : commentaires
\definecolor{qacodeannl}{HTML}{B26A00}     % code : annotations / javadoc

\color{qatextdark}

\setlength{\parskip}{4pt}
\setlength{\parindent}{0pt}

% ---------------------------------------------------------------
% TITRES DE SECTION -- texte blanc sur bande de couleur pleine largeur
% (gere aussi les variantes etoilees \section*)
% ---------------------------------------------------------------
\newcommand{\qasectionband}[2]{%
  \noindent\colorbox{#1}{%
    \parbox[c]{\dimexpr\linewidth-2\fboxsep\relax}{%
      \color{white}#2%
    }%
  }%
}
\newcommand{\secformat}[1]{\qasectionband{qaaccentdark}{\bfseries\Large\thesection\quad #1}}
\newcommand{\secformatnn}[1]{\qasectionband{qaaccentdark}{\bfseries\Large #1}}
\newcommand{\subsecformat}[1]{\qasectionband{qaaccent2dark}{\bfseries\large\thesubsection\quad #1}}
\newcommand{\subsecformatnn}[1]{\qasectionband{qaaccent2dark}{\bfseries\large #1}}

\titleformat{\section}{\normalfont}{}{0pt}{\secformat}
\titleformat{name=\section,numberless}{\normalfont}{}{0pt}{\secformatnn}
\titleformat{\subsection}{\normalfont}{}{0pt}{\subsecformat}
\titleformat{name=\subsection,numberless}{\normalfont}{}{0pt}{\subsecformatnn}
\titleformat{\subsubsection}{\normalfont\bfseries\color{qaaccentdark}}{\thesubsubsection}{0.6em}{}
\titlespacing*{\section}{0pt}{14pt}{8pt}
\titlespacing*{\subsection}{0pt}{10pt}{6pt}
\titlespacing*{\subsubsection}{0pt}{8pt}{4pt}

% ---------------------------------------------------------------
% LISTINGS -- theme clair (lstlisting reste utilisable tel quel)
% \lstcode : police + couleur de base, a reutiliser dans les surcharges
% ---------------------------------------------------------------
\newcommand{\lstcode}{\ttfamily\color{qatextdark}}

\lstdefinelanguage{TypeScript}{
  keywords={typeof, new, true, false, catch, function, return, null, switch, var,
    if, in, while, do, else, case, break, const, let, interface, class, extends,
    implements, enum, export, import, this, super, constructor, public, private,
    protected, static, async, await, throw, throws, try, finally, void},
  ndkeywords={boolean, any, number, string, int, long, double},
  sensitive=true,
  comment=[l]{//},
  morecomment=[s]{/*}{*/},
  morestring=[b]',
  morestring=[b]"
}

\lstdefinestyle{r307code}{
  basicstyle=\lstcode\footnotesize,
  keywordstyle=\color{qacodekwl}\bfseries,
  ndkeywordstyle=\color{qaaccent2dark}\bfseries,
  stringstyle=\color{qacodestrl},
  commentstyle=\color{qacodecoml}\itshape,
  morecomment=[s][\color{qacodeannl}]{/**}{*/},
  backgroundcolor=\color{qapanellight},
  rulecolor=\color{qalinelight},
  frame=single,
  framesep=4pt,
  xleftmargin=6pt,
  xrightmargin=6pt,
  numbers=left,
  numberstyle=\tiny\color{qamuteddark},
  numbersep=8pt,
  keepspaces=true,
  columns=fullflexible,
  showstringspaces=false,
  showspaces=false,
  showtabs=false,
  breaklines=true,
  breakatwhitespace=false,
  captionpos=b,
  tabsize=2,
  escapeinside={(*@}{@*)},
}
\lstset{style=r307code}

% ---------------------------------------------------------------
% BOITES : exercice, code, information
% ---------------------------------------------------------------
\tcbset{
  exobox/.style={colback=qapanellight, colframe=qaaccentdark, boxrule=0.6pt,
    arc=2pt, left=8pt, right=8pt, top=6pt, bottom=6pt,
    coltitle=white, colbacktitle=qaaccentdark, fonttitle=\bfseries,
    boxsep=2pt, breakable},
  infobox/.style={colback=white, colframe=qaaccent2dark, boxrule=0.6pt,
    leftrule=3pt, arc=0pt, left=8pt, right=8pt, top=5pt, bottom=5pt,
    coltitle=qaaccent2dark, colbacktitle=white, fonttitle=\bfseries,
    breakable},
  codebox/.style={colback=qatextdark, colframe=qalinelight, boxrule=0.4pt,
    arc=2pt, left=8pt, right=8pt, top=5pt, bottom=5pt, listing only,
    listing options={basicstyle=\ttfamily\small\color{white},
      breaklines=true, showstringspaces=false, columns=fullflexible,
      keepspaces=true, frame=none, numbers=none, backgroundcolor={}}, breakable}
}
\newtcolorbox[auto counter]{exobox}[2][]{exobox,
  title={Exercice~\thetcbcounter\ifstrempty{#2}{}{ --- #2}}, #1}
\newtcolorbox{infobox}[1][\`A retenir]{infobox, title={#1}}
\newtcblisting{codebox}{codebox}

% ---------------------------------------------------------------
% EN-TETE DE PAGE (fil d'Ariane + pagination)
% ---------------------------------------------------------------
\newcommand{\qaheadtitlevalue}{}
\newcommand{\qaheadtitle}[1]{\renewcommand{\qaheadtitlevalue}{#1}}

\pagestyle{fancy}
\fancyhf{}
\renewcommand{\headrulewidth}{0.4pt}
\fancyhead[L]{\small\color{qamuteddark}R3.07 --- \qaheadtitlevalue}
\fancyhead[R]{\small\color{qamuteddark}\thepage\,/\,\pageref{LastPage}}
\fancyfoot[C]{\footnotesize\color{qamuteddark}IUT de Vannes --- Universit\'e Bretagne Sud}
\fancypagestyle{plain}{\fancyhf{}\renewcommand{\headrulewidth}{0pt}%
  \fancyfoot[C]{\footnotesize\color{qamuteddark}IUT de Vannes --- Universit\'e Bretagne Sud}}

% ---------------------------------------------------------------
% BANDEAU DE TITRE (pas une page dediee, pour rester compact)
% \rtpheader{numero ou intitule court}{titre du TP}
% ---------------------------------------------------------------
\newcommand{\ryear}{2026/2027}
\newcommand{\rtpheader}[2]{%
  \thispagestyle{plain}%
  \begin{minipage}[t]{0.68\textwidth}
    {\small\color{qamuteddark}
      Universit\'e Bretagne Sud\\
      IUT de Vannes\\
      D\'epartement Informatique --- \ryear}
  \end{minipage}%
  \hfill
  \begin{minipage}[t]{0.28\textwidth}
    \raggedleft
    \includegraphics[width=2.6cm]{iut_logo.jpg}
  \end{minipage}

  \vspace{6mm}
  \begin{center}
    {\large\bfseries\color{qatextdark} BUT 2 --- R3.07. SQL dans un langage de programmation}\\[2mm]
    {\normalsize\color{qamuteddark} TD/TP #1}\\[1mm]
    {\Large\bfseries\color{qaaccentdark} #2}
  \end{center}
  \vspace{4mm}
  \qaheadtitle{#2}
}

\graphicspath{{../../latex/common/assets/}{./}}
\lstset{language=SQL}

\begin{document}
\rtpheader{6}{Proc\'edures stock\'ees et triggers}
\section*{Contexte général}

Lors des TP précédents, l’application bancaire a été développée en faisant évoluer
progressivement l’accès aux données :
\begin{itemize}
    \item accès direct via \textbf{JDBC} ;
    \item utilisation de \textbf{JPA} et de l’ORM ;
    \item structuration de l’accès aux données avec le \textbf{patron Repository}.
\end{itemize}

Dans tous ces cas, la logique métier était principalement implémentée
\textbf{côté application Java}.

Dans ce dernier TP, une partie de cette logique va être déplacée
\textbf{côté base de données}, à l’aide des fonctionnalités avancées des serveurs de bases de données (SGBD), dans notre cas PostgreSQL :
\begin{itemize}
    \item les \textbf{procédures stockées} ;
    \item et les \textbf{triggers}.
\end{itemize}

La base de données PostgreSQL sera exécutée dans un conteneur Docker, afin de garantir un environnement reproductible. Nous allons écrire ces deux fonctionnalités avec le langage PL/pgSQL (\textit{Procedural Language PostgreSQL}). Un mini-tutoriel est donné à la fin du sujet.

\bigskip

%-------------------------------------------------
\section*{Procédures stockées}

Une procédure stockée est un programme stocké et exécuté directement par le SGBD (dans notre cas, PostgreSQL).
Elle permet :
\begin{itemize}
    \item de regrouper plusieurs instructions SQL dans une seule unité logique ;
    \item de centraliser certaines règles métier ;
    \item de réduire les échanges entre l’application et la base de données ;
    \item d’améliorer la sécurité en contrôlant l’accès aux données.
\end{itemize}

Dans PostgreSQL, les procédures et fonctions peuvent être écrites en \texttt{PL/pgSQL}
et invoquées depuis une application Java via JDBC ou JPA.

\bigskip

%-------------------------------------------------
\section*{Triggers}

Un trigger est un mécanisme permettant d’exécuter automatiquement une action
lorsqu’un événement se produit sur une table (INSERT, UPDATE, DELETE, ...).

Les triggers sont principalement utilisés pour :
\begin{itemize}
    \item garantir des règles d’intégrité complexes ;
    \item tracer des modifications (journalisation) ;
    \item maintenir la cohérence des données ;
    \item appliquer des règles indépendantes de l’application.
\end{itemize}

Un trigger est associé à une table et repose sur une fonction déclenchée automatiquement.

\bigskip

%-------------------------------------------------
\section*{Schéma de la base de données}

On reprend les tables de l’application bancaire développée précédemment :
\begin{itemize}
    \item \texttt{Client(numClient, nomClient, prenomClient, ...)}
    \item \texttt{Compte(numCompte, solde, typeCompte)}
    \item \texttt{Compte\_Client(numClient, numCompte)}
    \item \texttt{Agent(numAgent, nomAgent, ...)}
    \item \texttt{Agence(numAgence, telAgence, ...)}
\end{itemize}

\bigskip

%-------------------------------------------------
\section*{Exercice 1 -- Procédures stockées}

On souhaite implémenter une opération de virement bancaire entre deux comptes.

\subsection*{Travail demandé}
\begin{enumerate}[label=\alph*)]
    \item Créer une procédure stockée \texttt{virement} prenant en paramètres :
    \begin{itemize}
        \item l’identifiant du compte source ;
        \item l’identifiant du compte destination ;
        \item le montant du virement.
    \end{itemize}

    \item La procédure devra :
    \begin{itemize}
        \item vérifier que le compte source dispose d’un solde suffisant ;
        \item débiter le compte source ;
        \item créditer le compte destination ;
        \item garantir l’atomicité de l’opération.
    \end{itemize}

    \item Tester la procédure directement dans PostgreSQL.

    \item Expliquer l’intérêt d’implémenter cette logique sous forme de procédure
    stockée plutôt que côté application Java.

    \item Que se passe-t-il si le compte \textbf{destination} n'existe pas ? Testez avec un numéro inexistant,
    puis regardez le solde du compte source. Corrigez la procédure (indice : la variable \texttt{FOUND}
    après un \texttt{UPDATE}).

    \item Appelez la procédure depuis votre application Java :
\begin{lstlisting}[language=Java]
try (PreparedStatement ps = c.prepareStatement("CALL virement(?, ?, CAST(? AS REAL))")) {
    ps.setInt(1, src);
    ps.setInt(2, dst);
    ps.setDouble(3, montant);
    ps.execute();
} catch (SQLException e) {
    System.err.println(e.getSQLState() + " : " + e.getMessage());
}
\end{lstlisting}
    Provoquez un solde insuffisant : quel \texttt{SQLState} reçoit Java ? Transformez cette erreur en
    exception métier \texttt{SoldeInsuffisantException}. Pourquoi le \texttt{CAST} est-il nécessaire
    (essayez sans) ?
\end{enumerate}
La correction est disponible sur Moodle. Pour bien en comprendre les détails, nous vous invitons à lire attentivement le sujet jusqu’à la fin.
\bigskip

%-------------------------------------------------
\section*{Exercice 2 -- Triggers}

On souhaite tracer automatiquement les opérations effectuées sur les comptes bancaires.

\subsection*{Travail demandé}

\begin{enumerate}[label=\alph*)]
    \item Créer une table \texttt{HistoriqueCompte} contenant :
    \begin{itemize}
        \item l’identifiant du compte ;
        \item l’ancienne valeur du solde ;
        \item la nouvelle valeur du solde ;
        \item la date de modification.
    \end{itemize}

    \item Créer une fonction de trigger permettant d’insérer automatiquement
    une ligne dans \texttt{HistoriqueCompte} lors d’une modification du solde d’un compte.

    \item Associer cette fonction à un trigger déclenché lors d’un \texttt{UPDATE}
    sur la table \texttt{Compte}.

    \item Tester le trigger à l’aide :
    \begin{itemize}
        \item d’une mise à jour SQL directe ;
        \item de la procédure stockée du premier exercice.
    \end{itemize}

    \item Expliquer l’intérêt d’un trigger par rapport à une journalisation
    réalisée côté application.
\end{enumerate}
La correction est disponible sur Moodle. Pour bien en comprendre les détails, nous vous invitons à lire attentivement le sujet jusqu’à la fin.
\bigskip

%-------------------------------------------------
\section*{Exercice 3 -- Curseurs : relevé de compte}

On souhaite produire le relevé d'un compte : ses opérations dans l'ordre chronologique, avec le
\textbf{solde cumulé} après chaque opération (un dépôt augmente le cumul, un retrait le diminue).

\subsection*{Travail demandé}
\begin{enumerate}[label=\alph*)]
    \item Écrire une fonction \texttt{releve\_compte(p\_num INT)} qui retourne une table
    \texttt{(date\_op DATE, type\_op VARCHAR, montant\_op REAL, cumul REAL)}. Elle doit parcourir les
    opérations du compte avec un \textbf{curseur explicite} (\texttt{OPEN}, \texttt{FETCH}, \texttt{CLOSE})
    et renvoyer chaque ligne avec \texttt{RETURN NEXT}.
    \item Lever une exception si le compte n'existe pas.
    \item Tester : \texttt{SELECT * FROM releve\_compte(1);}
    \item Obtenir le même résultat avec \textbf{une seule requête SQL}, sans curseur, grâce à une fonction de fenêtre :
    \texttt{sum(\ldots) OVER (ORDER BY dateOperation, numOperation)}. Quelle version préférez-vous, et pourquoi ?
    \item Depuis Java, afficher le relevé avec un \texttt{PreparedStatement} sur
    \texttt{SELECT * FROM releve\_compte(?)}.
\end{enumerate}

\bigskip

\section*{Rendu attendu}

\begin{itemize}
    \item Les scripts SQL (procédures, triggers, tests) sont donnés sur Moodle. Pas besoin de les rendre ;
    \item Résultats des tests (captures d'écran) ;
    \item Réponses argumentées aux questions dans un fichier Markdown.
\end{itemize}

\section*{Exemples de procédures stockées et triggers}

\subsection*{1. Exemples de procédures stockées (logique métier explicite)}

Les procédures stockées sont idéales pour encapsuler des règles métier critiques, transactionnelles et réutilisables par toutes les applications clientes.

\begin{enumerate}

\item \textbf{Clôture d'un compte bancaire}

\textbf{Procédure :} \texttt{cloturer\_compte(numCompte)}

\begin{itemize}
    \item Vérifie que le solde est nul.
    \item Empêche la clôture si le compte est débiteur.
    \item Supprime ou archive le compte.
\end{itemize}

\textbf{Pourquoi une procédure ?}
Règle métier critique, transactionnelle, commune à toutes les applications clientes.

\item \textbf{Ouverture d'un compte}

\textbf{Procédure :} \texttt{ouvrir\_compte(typeCompte, soldeInitial)}

\begin{itemize}
    \item Crée un nouveau compte.
    \item Applique des règles selon le type de compte (minimum initial, découvert autorisé).
    \item Initialise correctement les données.
\end{itemize}

\textbf{Pourquoi une procédure ?}
Centralise les règles d’ouverture indépendamment du code Java.

\item \textbf{Changement de type de compte}

\textbf{Procédure :} \texttt{changer\_type\_compte(numCompte, nouveauType)}

\begin{itemize}
    \item Vérifie la compatibilité du type avec le solde actuel.
    \item Empêche certains changements (ex : compte épargne → compte courant).
\end{itemize}

\textbf{Pourquoi une procédure ?}
Règles complexes dépendantes des données existantes.

\item \textbf{Calcul et versement des intérêts}

\textbf{Procédure :} \texttt{verser\_interets()}

\begin{itemize}
    \item Parcourt tous les comptes épargne.
    \item Calcule les intérêts.
    \item Met à jour les soldes.
\end{itemize}

\textbf{Pourquoi une procédure ?}
Traitement batch, exécuté périodiquement côté base (cron, tâche planifiée).

\item \textbf{Virement groupé}

\textbf{Procédure :} \texttt{virement\_multiple(listeComptes, montant)}

\begin{itemize}
    \item Effectue plusieurs virements dans une seule transaction.
    \item Annule tout si une seule opération échoue.
\end{itemize}

\textbf{Pourquoi une procédure ?}
Garantie d’atomicité sur plusieurs opérations.

\end{enumerate}

\subsection*{2. Exemples de triggers (règles automatiques et transverses)}

Les triggers sont utilisés pour appliquer des règles automatiques, transverses et non contournables.

\begin{enumerate}

\item \textbf{Interdiction de solde négatif}

\textbf{Trigger :} \texttt{before\_update\_solde}

\begin{itemize}
    \item Vérifie qu’un solde ne devient pas négatif.
    \item Bloque l’opération automatiquement.
\end{itemize}

\textbf{Pourquoi un trigger ?}
Règle d’intégrité qui doit s’appliquer même en cas d’UPDATE SQL direct.

\item \textbf{Journalisation des opérations}

\textbf{Trigger :} \texttt{after\_update\_compte}

\begin{itemize}
    \item Enregistre toute modification de solde.
    \item Sert à l’audit et à la traçabilité.
\end{itemize}

\textbf{Pourquoi un trigger ?}
Indépendant de l’application, impossible à contourner.

\item \textbf{Détection d’opérations suspectes}

\textbf{Trigger :} \texttt{after\_update\_solde}

\begin{itemize}
    \item Détecte des montants inhabituels.
    \item Insère une alerte dans une table dédiée.
\end{itemize}

\textbf{Pourquoi un trigger ?}
Surveillance automatique et continue.

\item \textbf{Mise à jour d’un solde global}

\textbf{Trigger :} \texttt{after\_update\_compte}

\begin{itemize}
    \item Met à jour le solde total d’une agence.
    \item Maintient des données agrégées cohérentes.
\end{itemize}

\textbf{Pourquoi un trigger ?}
Évite des calculs coûteux répétés côté application.

\item \textbf{Historisation des comptes supprimés}

\textbf{Trigger :} \texttt{before\_delete\_compte}

\begin{itemize}
    \item Archive les informations du compte supprimé.
    \item Permet une récupération ou un audit ultérieur.
\end{itemize}

\textbf{Pourquoi un trigger ?}
Traçabilité forte, aucune suppression silencieuse.

\end{enumerate}

\subsection*{Message clé}

\begin{itemize}
\item \textbf{Les procédures stockées} sont adaptées aux règles métier explicites et transactionnelles ;
\item \textbf{Les triggers} sont adaptés aux règles automatiques, transverses et non contournables.
\end{itemize}
        

\section*{Mini-tutoriel : bases du PL/pgSQL}

PL/pgSQL (Procedural Language PostgreSQL) est un langage procédural permettant
d’écrire des programmes exécutés directement par PostgreSQL.
Il est utilisé pour définir des \textbf{procédures stockées} et des \textbf{triggers}.

Ce tutoriel présente uniquement les éléments nécessaires à la réalisation de ce TP.

\bigskip

\subsection*{Structure générale d’une procédure}

Une procédure stockée PostgreSQL se définit avec l’instruction
\texttt{CREATE PROCEDURE} et contient un bloc \texttt{BEGIN ... END}.

\begin{verbatim}
CREATE OR REPLACE PROCEDURE nom_procedure(paramètres)
LANGUAGE plpgsql
AS $$
BEGIN
    -- instructions SQL
END;
$$;
\end{verbatim}

Une procédure peut être appelée avec l’instruction \texttt{CALL}.

\begin{verbatim}
CALL nom_procedure(...);
\end{verbatim}

\bigskip

\subsection*{Déclaration de variables}

Les variables sont déclarées dans une section \texttt{DECLARE}.

\begin{verbatim}
DECLARE
    solde_source REAL;
\end{verbatim}

Elles peuvent ensuite être utilisées dans le corps de la procédure.

\bigskip

\subsection*{Récupération d’un résultat dans une variable}

L’instruction \texttt{SELECT ... INTO} permet de stocker le résultat
d’une requête SQL dans une variable.

\begin{verbatim}
SELECT solde INTO solde_source
FROM Compte
WHERE numCompte = 1001;
\end{verbatim}

\bigskip

\subsection*{Conditions}

Les structures conditionnelles utilisent \texttt{IF / THEN / END IF}.

\begin{verbatim}
IF solde_source < 0 THEN
    RAISE EXCEPTION 'Solde négatif';
END IF;
\end{verbatim}

\bigskip

\subsection*{Gestion des erreurs}

L’instruction \texttt{RAISE EXCEPTION} permet d’arrêter l’exécution
de la procédure et d’afficher un message d’erreur.

\begin{verbatim}
RAISE EXCEPTION 'Erreur : solde insuffisant';
\end{verbatim}

\bigskip

\subsection*{Transactions}

Dans PostgreSQL, une procédure stockée s’exécute par défaut
dans une transaction.  
Si une exception est levée, toutes les modifications sont annulées.

Cela garantit l’atomicité des opérations sensibles comme un virement bancaire.

\bigskip

\subsection*{Curseurs}

Un curseur parcourt le résultat d'une requête \textbf{ligne par ligne}. On le déclare, on l'ouvre,
on lit chaque ligne avec \texttt{FETCH}, et on s'arrête quand \texttt{FOUND} devient faux.

\begin{verbatim}
DECLARE
    cur CURSOR FOR SELECT numCompte, solde FROM Compte ORDER BY numCompte;
    ligne RECORD;
BEGIN
    OPEN cur;
    LOOP
        FETCH cur INTO ligne;
        EXIT WHEN NOT FOUND;
        RAISE NOTICE 'Compte % : %', ligne.numCompte, ligne.solde;
    END LOOP;
    CLOSE cur;
END;
\end{verbatim}

La forme \texttt{FOR ligne IN SELECT \ldots{} LOOP \ldots{} END LOOP;} utilise un curseur implicite.
Une fonction qui renvoie plusieurs lignes se déclare \texttt{RETURNS TABLE(\ldots)} et renvoie chaque
ligne avec \texttt{RETURN NEXT}. Attention : un traitement ligne par ligne est souvent plus lent qu'une
requête SQL ensembliste équivalente.

\bigskip

\subsection*{Fonctions et triggers}

Un trigger repose sur une \textbf{fonction} qui retourne \texttt{TRIGGER}.

\begin{verbatim}
CREATE OR REPLACE FUNCTION nom_fonction()
RETURNS TRIGGER
LANGUAGE plpgsql
AS $$
BEGIN
    -- traitement automatique
    RETURN NEW;
END;
$$;
\end{verbatim}

Dans une fonction de trigger :
\begin{itemize}
    \item \texttt{OLD} représente l’ancienne ligne dans la table ;
    \item \texttt{NEW} représente la nouvelle ligne.
\end{itemize}

\bigskip

\subsection*{Création d’un trigger}

Un trigger est associé à une table et à un événement.

\begin{verbatim}
CREATE TRIGGER nom_trigger
AFTER UPDATE
ON Compte
FOR EACH ROW
EXECUTE FUNCTION nom_fonction();
\end{verbatim}

Le code de la fonction est exécuté automatiquement
lors de l’événement concerné.


\end{document}
